\chapter{Essential Tremor}
\index{Essential Tremor}

\section*{Definition and Overview}
Essential tremor is a common neurological condition that causes a rhythmic, involuntary shaking, most often of the hands and arms. It is an "action tremor," meaning it is most noticeable when the person is using the hands, such as holding a cup, writing, or doing fine tasks, rather than when the hands are at rest. Essential tremor is not dangerous and does not shorten life, but it can be embarrassing and interfere with daily activities. Although there is no cure, medicines and, in severe cases, procedures can significantly reduce the tremor.

\section*{Epidemiology}
Essential tremor is one of the most common movement disorders, affecting about 1 per cent of the population overall and a higher proportion of older people, with some studies finding it in up to 5 per cent of those over 60. It affects men and women equally and can begin at any age, but it becomes more common with advancing years. Despite its frequency, it is often not diagnosed or treated, because many people accept the shaking as a normal part of ageing.

\section*{Aetiology and Pathophysiology}
The cause of essential tremor is not fully understood.
\begin{itemize}
    \item \textbf{Genetics:} essential tremor runs in families, and several genes have been linked to it
    \item \textbf{Brain circuits:} research points to abnormal activity in the circuits of the brain that control movement, particularly the cerebellum
    \item \textbf{Mechanism:} the tremor results from abnormal rhythmic signals in the motor pathways
    \item \textbf{Triggers and aggravators:} stress, anxiety, fatigue, caffeine, and some medicines worsen the tremor
    \item \textbf{Alcohol:} alcohol often temporarily reduces essential tremor, which is a diagnostic clue, though it is not a treatment
    \item \textbf{Progression:} the tremor tends to worsen gradually with age
\end{itemize}

\section*{Clinical Features}
Essential tremor has a characteristic pattern.
\begin{itemize}
    \item Rhythmic shaking of the hands and arms, usually on both sides
    \item The tremor appears with action, such as holding a cup, writing, or pointing
    \item It may affect the head, with a yes-yes or no-no nodding
    \item The voice can be shaky
    \item The legs and trunk are less commonly affected
    \item Tremor worsens with stress, fatigue, and caffeine
    \item Difficulty with fine tasks: writing, eating, drinking, and dressing
    \item A rest tremor is not typical, which helps distinguish it from Parkinson's disease
\end{itemize}

\section*{Differential Diagnosis}
Other causes of tremor must be distinguished.
\begin{itemize}
    \item Parkinson's disease, which causes a resting tremor, slowness, and stiffness
    \item Tremor from medicines, including some asthma and psychiatric medicines
    \item Tremor from thyroid disease, low blood sugar, or liver disease
    \item Cerebellar disorders, with incoordination
    \item Physiologic tremor, which is normal and mild, worsened by caffeine or fatigue
    \item Functional tremor
\end{itemize}

\section*{When to Seek Medical Care}
People with a tremor that interferes with daily life, or that is new, progressive, or accompanied by other symptoms, should be assessed by a GP and referred to a neurologist.

\textbf{Contact your local emergency services for:}
\begin{itemize}
    \item Sudden onset of tremor with weakness, facial drooping, or difficulty speaking, which may indicate stroke
    \item A tremor with sudden confusion
    \item Collapse or loss of consciousness
\end{itemize}

\section*{Diagnosis and Investigations}
Diagnosis is clinical, based on the pattern of the tremor.
\begin{itemize}
    \item \textbf{History and examination:} the type, timing, and triggers of the tremor, and a family history
    \item \textbf{Observation of the tremor:} during writing, holding objects, and at rest
    \item \textbf{Blood tests:} thyroid function and other checks to exclude causes
    \item \textbf{Exclusion of Parkinson's disease:} assessment for stiffness, slowness, and rest tremor
    \item \textbf{Specialist referral:} to a neurologist for confirmation and management
\end{itemize}

\section*{Management}
Management aims to reduce the impact of the tremor.
\begin{itemize}
    \item \textbf{Reassurance:} understanding that essential tremor is not dangerous
    \item \textbf{Reducing triggers:} limiting caffeine, managing stress, and ensuring adequate sleep
    \item \textbf{Medicines:} beta-blockers such as propranolol and other agents reduce tremor in many people
    \item \textbf{Occupational therapy:} weighted utensils, adapted tools, and strategies for daily tasks
    \item \textbf{Aids:} heavier cups, writing aids, and assistive devices
    \item \textbf{Botulinum toxin injections:} for head and voice tremor
    \item \textbf{Procedures:} deep brain stimulation or focused ultrasound for severe, disabling tremor
\end{itemize}

\section*{Complications}
\begin{itemize}
    \item Difficulty with eating, drinking, writing, and personal care
    \item Embarrassment and social withdrawal
    \item Reduced independence and quality of life
    \item Anxiety and depression
    \item Side effects of medicines
\end{itemize}

\section*{Prognosis and Outcomes}
Essential tremor is a chronic condition that usually worsens slowly, but it does not affect life expectancy. Many people manage well with medicines and simple aids, and those with severe tremor can benefit greatly from modern procedures. The condition is often mild, and treatment is tailored to the degree to which the tremor affects daily life.

\section*{Prevention and Self-Care}
\begin{itemize}
    \item Limit caffeine and alcohol
    \item Manage stress and fatigue
    \item Use weighted or adapted utensils and cups
    \item Practise techniques such as holding objects close to the body
    \item Seek occupational therapy advice
    \item Take medicines as prescribed
    \item Discuss severe tremor with a neurologist
\end{itemize}

\section*{Sources and Further Reading}
The following reputable sources provide further information for healthcare students and patients seeking reliable, current advice about essential tremor.
\begin{itemize}
    \item Essential Tremor Australia and International Essential Tremor Foundation. \textit{Information and support.} https://www.essentialtremor.org/
    \item Healthdirect Australia. \textit{Tremor.} https://www.healthdirect.gov.au/tremor
    \item NHS. \textit{Essential tremor.} https://www.nhs.uk/conditions/essential-tremor/
    \item Mayo Clinic. \textit{Essential tremor.} https://www.mayoclinic.org/diseases-conditions/essential-tremor/
    \item MedlinePlus. \textit{Essential tremor.} https://medlineplus.gov/essentialtremor.html
    \item MSD Manual. \textit{Tremor.} https://www.msdmanuals.com/
\end{itemize}
